\clearpage
\VerdierPhysicalPageAnchor{0154}%
\begingroup
\small
\setlength{\parskip}{0pt}%
\setlength{\emergencystretch}{1.5em}%

\VerdierSemanticAnchor{ch3-3-2}
\subsection{Propriétés de la catégorie
  \texorpdfstring{$\mathrm A(\mathcal D)$}{A(D)}
  (\texorpdfstring{$\mathcal D$}{D} catégorie triangulée).}

\VerdierSemanticAnchor{ch3-3-2-0}
\noindent\textbf{3.2.0.}\quad
Soient $\mathcal D$ et $\mathcal D'$ deux catégories triangulées. Désignons
par $\mathcal E x(\mathcal D,\mathcal D')$ la catégorie dont les objets sont
les foncteurs exacts de $\mathcal D$ dans $\mathcal D'$ et dont les
morphismes sont les morphismes de foncteurs exacts. Désignons de même par
$\mathbb Z\text{-}\operatorname{Add}(\mathrm A(\mathcal D),
\mathrm A(\mathcal D'))$ la catégorie dont les objets sont les
$\mathbb Z$-foncteurs additifs de $\mathrm A(\mathcal D)$ dans
$\mathrm A(\mathcal D')$ et dont les morphismes sont les morphismes de
$\mathbb Z$-foncteurs. Il résulte de ce qui précède qu'on a défini un
foncteur

\VerdierSemanticAnchor{ch3-3-2-0-1}
\[
\mathrm A_{\mathcal D,\mathcal D'}\colon
\mathcal E x(\mathcal D,\mathcal D')
\longrightarrow
\mathbb Z\text{-}\operatorname{Add}
(\mathrm A(\mathcal D),\mathrm A(\mathcal D')),
\tag{3.2.0.1}
\]
qui associe à tout foncteur exact $F$ le foncteur $\mathrm A(F)$.

\VerdierSemanticAnchor{ch3-3-2-1}
\medskip
\noindent\textbf{Théorème 3.2.1.}\begingroup\renewcommand{\thefootnote}{(*)}%
\footnote{Ce théorème m'a été communiqué par Freyd [17].}\endgroup\quad
\textit{a) Si $\mathcal D$ est une catégorie triangulée, la catégorie
$\mathrm A(\mathcal D)$ est abélienne. Si
$F\colon\mathcal D\to\mathcal D'$ est un foncteur exact entre deux
catégories triangulées, le foncteur $\mathrm A(F)$ (3.1.4.1) est exact. Le
foncteur $\mathrm A_{\mathcal D,\mathcal D'}$ (3.2.0.1) est pleinement fidèle
et injectif sur les objets.}

\VerdierSemanticAnchor{ch3-3-2-1-b}
\textit{b) Les objets de $\mathcal D$ sont injectifs et projectifs dans
$\mathrm A(\mathcal D)$. Tout objet de $\mathrm A(\mathcal D)$ est isomorphe
à un objet quotient (resp. à un sous-objet) d'un objet de $\mathcal D$.}

\VerdierSemanticAnchor{ch3-3-2-1-c}
\textit{c) Soit $\mathcal A$ une catégorie abélienne. Désignons par
$\mathcal E x(\mathrm A(\mathcal D),\mathcal A)$ la catégorie dont les objets
sont les foncteurs exacts de $\mathrm A(\mathcal D)$ dans $\mathcal A$ et les
morphismes, les morphismes de foncteurs. De même, désignons par
$\mathcal H\!om(\mathcal D,\mathcal A)$ la catégorie dont les objets sont les
foncteurs de $\mathcal D$ dans $\mathcal A$ et les morphismes, les morphismes
de foncteurs. La composition avec le foncteur
$\mathcal H\colon\mathcal D\to\mathrm A(\mathcal D)$ définit un foncteur}
\[
\mathcal E x(\mathrm A(\mathcal D),\mathcal A)
\longrightarrow
\mathcal H\!om(\mathcal D,\mathcal A).
\]
\textit{Ce foncteur est pleinement fidèle. Son image essentielle est constituée
des foncteurs cohomologiques de $\mathcal D$ dans $\mathcal A$ (1.1.5). En
particulier, $\mathcal H\colon\mathcal D\to\mathrm A(\mathcal D)$ est un
foncteur cohomologique.}

\VerdierSemanticAnchor{ch3-3-2-1-proof-plan}
La démonstration de ce théorème nous occupera jusqu'à l'alinéa (3.2.9). Nous
démontrerons d'abord trois propositions intermédiaires.

\VerdierSemanticAnchor{ch3-3-2-2}
\medskip
\noindent\textbf{3.2.2.}\quad
Soit $U$ un univers tel que $\mathcal D$ soit $U$-petite, c'est-à-dire tel que
$\mathcal D$ soit un élément de $U$. Soit
$\mathcal D^\wedge$ la catégorie des foncteurs additifs contravariants sur
$\mathcal D$ à valeurs dans la catégorie des groupes abéliens $U$-petits. La catégorie
$\mathcal D^\wedge$ est abélienne et le foncteur canonique
$h\colon\mathcal D\to\mathcal D^\wedge$, qui associe à tout objet de
$\mathcal D$ le foncteur
qu'il représente, est pleinement fidèle. Le foncteur $h$ transforme tout
objet de $\mathcal D$ en objet projectif de $\mathcal D^\wedge$.
\endgroup
